\documentclass[a4paper,11pt]{article}
\usepackage[italian]{babel}
% Shared preamble for the RESEARCH blog papers and their Italian explainers.
% Usage (after \documentclass[a4paper,11pt]{article} and \usepackage[<lang>]{babel}):
%   % Shared preamble for the RESEARCH blog papers and their Italian explainers.
% Usage (after \documentclass[a4paper,11pt]{article} and \usepackage[<lang>]{babel}):
%   % Shared preamble for the RESEARCH blog papers and their Italian explainers.
% Usage (after \documentclass[a4paper,11pt]{article} and \usepackage[<lang>]{babel}):
%   \input{../stile}
% Compile with: latexmk -xelatex -interaction=nonstopmode -halt-on-error <file>.tex
\usepackage[utf8]{inputenc}
\usepackage[T1]{fontenc}
\usepackage{lmodern}
\usepackage{csquotes}
\usepackage{microtype}
\usepackage[a4paper,margin=2.5cm]{geometry}
\usepackage{amsmath,amssymb}
\usepackage{booktabs}
\usepackage{array}
\usepackage{float}
\usepackage[skip=0.5\baselineskip plus 2pt]{parskip}
\usepackage{setspace}
\setstretch{1.12}
\usepackage{pgfplots}
\pgfplotsset{compat=1.18}
\usepgfplotslibrary{groupplots}
\usepackage{caption}
\captionsetup{font=small,labelfont=sc}
\usepackage[hidelinks]{hyperref}
\widowpenalty=10000
\clubpenalty=10000
\setlength{\emergencystretch}{2em}
\newcommand{\bp}{\,\mathrm{bp}}

% Compile with: latexmk -xelatex -interaction=nonstopmode -halt-on-error <file>.tex
\usepackage[utf8]{inputenc}
\usepackage[T1]{fontenc}
\usepackage{lmodern}
\usepackage{csquotes}
\usepackage{microtype}
\usepackage[a4paper,margin=2.5cm]{geometry}
\usepackage{amsmath,amssymb}
\usepackage{booktabs}
\usepackage{array}
\usepackage{float}
\usepackage[skip=0.5\baselineskip plus 2pt]{parskip}
\usepackage{setspace}
\setstretch{1.12}
\usepackage{pgfplots}
\pgfplotsset{compat=1.18}
\usepgfplotslibrary{groupplots}
\usepackage{caption}
\captionsetup{font=small,labelfont=sc}
\usepackage[hidelinks]{hyperref}
\widowpenalty=10000
\clubpenalty=10000
\setlength{\emergencystretch}{2em}
\newcommand{\bp}{\,\mathrm{bp}}

% Compile with: latexmk -xelatex -interaction=nonstopmode -halt-on-error <file>.tex
\usepackage[utf8]{inputenc}
\usepackage[T1]{fontenc}
\usepackage{lmodern}
\usepackage{csquotes}
\usepackage{microtype}
\usepackage[a4paper,margin=2.5cm]{geometry}
\usepackage{amsmath,amssymb}
\usepackage{booktabs}
\usepackage{array}
\usepackage{float}
\usepackage[skip=0.5\baselineskip plus 2pt]{parskip}
\usepackage{setspace}
\setstretch{1.12}
\usepackage{pgfplots}
\pgfplotsset{compat=1.18}
\usepgfplotslibrary{groupplots}
\usepackage{caption}
\captionsetup{font=small,labelfont=sc}
\usepackage[hidelinks]{hyperref}
\widowpenalty=10000
\clubpenalty=10000
\setlength{\emergencystretch}{2em}
\newcommand{\bp}{\,\mathrm{bp}}


\title{Un segnale che funziona e non paga\\[2pt]
\large Spiegazione semplice}
\author{Tommaso Zeri}
\date{Novembre 2026}

\begin{document}
\maketitle

\section*{Di che cosa parliamo}

Su Binance il bitcoin si compra in due modi. Il primo è lo \emph{spot}: compri il bitcoin e lo possiedi. Il secondo è il \emph{perpetual}, un contratto che segue il prezzo del bitcoin senza che tu lo possieda davvero. I due prezzi sono quasi uguali, ma non del tutto. La loro differenza si chiama \emph{basis} e oscilla di circa un basis point, cioè un centesimo di punto percentuale. Su un bitcoin da centomila dollari un basis point vale dieci dollari.

Quando il prezzo del perpetual si muove, quello dello spot di solito lo segue entro pochi millisecondi o al massimo un secondo. Per un istante quindi la basis si allarga, e poi si richiude. Il modello della tesi (un processo di Hawkes) descrive proprio questo: ogni movimento di prezzo rende più probabili altri movimenti, sullo stesso mercato o sull'altro. Una buona analogia è una sala dove qualcuno ride e altri ridono subito dopo. Se conosci chi ha riso e quando, puoi prevedere chi riderà nei prossimi secondi.

Da questo modello si ricava una previsione: \emph{nei prossimi dieci secondi la basis salirà o scenderà?} La domanda dell'articolo è se con questa previsione si guadagna.

\section*{Come si controlla senza barare}

Il test si chiama \emph{backtest}: si prende la storia passata e si simula che cosa sarebbe successo seguendo la regola. Il pericolo principale è il \emph{look-ahead}, cioè usare per decidere un'informazione che in quel momento non avresti avuto. Sarebbe come scommettere su una partita conoscendo già il risultato.

Per evitarlo ogni giornata è divisa in sei blocchi di quattro ore. Il modello impara dal blocco precedente e viene usato sul blocco successivo, che non ha mai visto. Ogni decisione usa soltanto i movimenti di prezzo già avvenuti. Le soglie della regola sono fissate prima, su una griglia, e non vengono ritoccate per far tornare i conti. Il test copre dieci giornate di bitcoin tra dicembre 2022 e settembre 2026 e due di ether.

\section*{Primo risultato: la previsione funziona}

Se si potesse comprare e vendere al prezzo medio, il \emph{mid}, la regola guadagnerebbe tra $0{,}4$ e $1{,}3$ basis point per operazione. Ma al prezzo medio non si compra. Chi vuole comprare subito paga il prezzo più alto esposto (l'\emph{ask}), chi vuole vendere subito incassa il più basso (il \emph{bid}). La differenza tra i due è lo \emph{spread}. Pagando lo spread su entrambi i mercati, il guadagno scende a $0{,}1$--$0{,}7$ basis point per operazione, $0{,}44$ in media sui giorni di bitcoin. È positivo in tutti i dodici giorni del test e resta circa uguale se gli ordini arrivano con un ritardo (\emph{latenza}) fino a dieci millisecondi.

Questi numeri sono \emph{lordi}: non contengono ancora le commissioni dell'exchange.

\section*{Perché lo spread costa più del previsto}

In media, attraversare lo spread dei due mercati costa tra $0{,}08$ e $0{,}26$ basis point. Il segnale però ne perde circa $0{,}40$ passando dal prezzo medio al prezzo reale. Il motivo è il momento in cui il segnale decide. Decide subito dopo un movimento di prezzo, e proprio in quell'istante lo spread del mercato che si è appena mosso è largo due o tre volte il normale. È come voler comprare un ombrello quando comincia a piovere: il prezzo è più alto proprio quando ti serve.

\section*{La regola ingenua e la selezione avversa}

Una regola più semplice dice: \enquote{il perpetual si è mosso, lo spot non ancora, compro lo spot al vecchio prezzo prima che si aggiorni}. Al prezzo medio questa regola sembra la migliore, con oltre un basis point per operazione. Al prezzo reale perde in dieci giorni su dodici.

Il fenomeno si chiama \emph{selezione avversa}. Pensa a un banco del mercato con un cartellino del prezzo vecchio. Se tutti sanno che il prezzo vero è già cambiato, il primo ad arrivare compra a quel prezzo vecchio, e il primo non sei tu: è chi ha il computer più vicino e più veloce. A te restano le occasioni che nessuno voleva, cioè quelle che fanno perdere. Più un segnale è pubblico e ovvio, più è probabile che qualcun altro lo sfrutti prima.

\section*{Le commissioni mangiano tutto}

Ogni operazione sulla basis richiede quattro scambi: si apre e si chiude una posizione su ciascuno dei due mercati. Binance fa pagare una commissione su ogni scambio. Il costo di un'operazione completa va da circa $2{,}2$ basis point per il cliente più grande che lascia ordini nel libro, a $8$ per lo stesso cliente se prende gli ordini degli altri, fino a $30$ per un conto normale. Il conto in tasca si fa così:
\[
\text{guadagno netto} = \text{guadagno lordo} - \text{commissioni} \approx 0{,}44 - 2{,}2 < 0 .
\]
Nel caso più favorevole il guadagno lordo è circa un quarto della commissione. Per un conto normale la commissione è più di dieci volte il guadagno. Nel test la regola indovina la direzione circa sette volte su dieci. Somiglia a una scommessa che vinci spesso, ma in cui ogni puntata costa più di quanto vinci: più giochi, più perdi.

\section*{E se si aspettasse invece di inseguire?}

Chi lascia un ordine fermo nel libro (il \emph{maker}) non paga lo spread e paga commissioni più basse. Sembrerebbe la soluzione. Il problema è \emph{quando} l'ordine viene eseguito. Per stimarlo la tesi usa una regola prudente: un ordine fermo conta come eseguito solo quando il prezzo del mercato lo supera.

Il risultato è istruttivo. Le poche operazioni in cui entrambi gli ordini vengono eseguiti guadagnano circa un basis point. Ma sono una minoranza, e si riconoscono solo a posteriori. Nella maggior parte dei casi i due mercati si muovono insieme, viene eseguito solo l'ordine dalla parte sbagliata e si resta con mezza posizione mentre il prezzo va contro. Queste operazioni perdono in media $3{,}4$ basis point. Dopo le commissioni nessuna variante chiude in attivo in nessuno dei dodici giorni.

\section*{Glossario}

\begin{description}
\item[basis] differenza tra il prezzo del perpetual e quello dello spot.
\item[basis point (bp)] un centesimo di punto percentuale, $0{,}01\%$.
\item[bid / ask] il prezzo più alto a cui qualcuno compra e il più basso a cui qualcuno vende in quel momento.
\item[mid] il punto medio tra bid e ask; un prezzo di riferimento, non un prezzo a cui si scambia.
\item[spread] la distanza tra ask e bid; è il costo di comprare e rivendere subito.
\item[touch] i prezzi migliori esposti nel libro, cioè il bid e l'ask; \enquote{eseguire al touch} vuol dire pagare lo spread.
\item[maker / taker] il maker lascia un ordine fermo e aspetta; il taker accetta subito un ordine già esposto.
\item[latenza] il ritardo tra la decisione e l'arrivo dell'ordine al mercato.
\item[backtest] simulazione di una strategia sui dati del passato.
\item[look-ahead] errore di un backtest che usa informazioni non ancora disponibili al momento della decisione.
\item[selezione avversa] la tendenza a essere eseguiti proprio negli scambi che fanno perdere, perché chi sta dall'altra parte sa qualcosa in più o arriva prima.
\item[processo di Hawkes] modello in cui ogni evento rende più probabili altri eventi nell'immediato futuro.
\end{description}

\section*{Cosa portare a casa}

Una previsione può essere giusta in modo misurabile e ripetibile e non valere niente in pratica. Qui il modello indovina la direzione della basis più spesso di quanto farebbe il caso, ogni giorno, ma il vantaggio è di mezzo basis point, mentre aprire e chiudere un'operazione costa da due a trenta basis point di commissioni. Il segnale scatta proprio quando lo spread è più caro, e la regola più ovvia viene battuta da chi è più veloce. Prima di credere a un backtest conviene chiedersi a quale prezzo si sarebbe davvero scambiato e quanto sarebbe costato farlo.

Questo testo è una spiegazione divulgativa di un lavoro di ricerca e non è una raccomandazione di investimento. È tratto dalla tesi magistrale dell'autore.
\end{document}
